% RESEARCH_DRAFT_SHA256: ef498131e9a981419b05c824c6dcba25d81ab05144698a7e7e68c70b02b1bd9f
% Layout and source-reference wrappers expanded for the web renderer.
\documentclass[11pt]{article}
\usepackage[margin=25mm]{geometry}
\usepackage{amsthm}
\usepackage{fontspec,unicode-math}
\setmainfont{Latin Modern Roman}
\setmathfont{Latin Modern Math}
\usepackage{amsmath,amssymb,mathtools,xurl,hyperref}
\hypersetup{colorlinks=true,urlcolor=blue,linkcolor=blue}
\setcounter{secnumdepth}{0}
\setlength{\parindent}{0pt}
\setlength{\parskip}{5pt}
\setlength{\emergencystretch}{3em}
\newcommand{\sref}[2]{\hyperlink{source:#1:#2}{[S#2]}}
\newcommand{\cataloguescope}[1]{\paragraph{Source record.}#1}
\newcommand{\floor}[1]{\left\lfloor#1\right\rfloor}
\newtheorem{proposition}{Proposition}
\newtheorem{lemma}[proposition]{Lemma}
\begin{document}
\paragraph{Research manuscript.} OpenAI \texttt{gpt-6-astra}, reasoning \texttt{ultra}; 11 October 2026.
\paragraph{Preserved source body.} The complete catalogue mathematical text and sources follow with layout and source links expanded, including source qualifications and its historical preparation notes. New research begins at the section headed Research attempt.

\section{Mixing of Random Colorings on General Graphs}\label{statement}

\subsection{Definitions and mathematical statement}

\paragraph{Shared notation.}Standard number systems are $\mathbb N,\mathbb Z,\mathbb Q,\mathbb R,\mathbb C$. Specialized notation and index ranges follow the cited source and the definitions below.


A proper $q$-coloring assigns one of $q$ colors to each vertex with unequal colors on adjacent vertices. Spatial mixing means boundary-condition influence decays with distance; temporal mixing means a sampling chain approaches its stationary coloring law.
\paragraph{Record scope and qualifications.}
The catalogue title is retained as metadata. Named certificates, contracts, bounds or frameworks in that title are not asserted as results of the original source extract.
The number of colors, degree range, coloring measure, dynamics, and desired rate are not specified in the supplied question.

Source-record limitation: The graph/color range, dynamics, stationary law, and requested mixing rate are not specified.
\paragraph{Statement recorded in the supplied source.}
Establish spatial and temporal mixing properties for random proper colorings on general graphs. \href{https://doi.org/10.4171/owr/2024/27}{[S2]}
\paragraph{Source status.}
The supplied catalogue labels this source record open; that metadata was checked against the supplied snapshot on 2026-10-11. The label is preserved for provenance and is not an independent verification of present open status or mathematical correctness. The catalogue literature assessment is dated 2026-08-22.
\subsection{Short English statement}
Prove spatial and temporal mixing for random proper colorings on the stated graph classes.
\subsection{Sources}
\begin{itemize}
\item[S1] Ulam AI and the UnsolvedMath Contributors, UnsolvedMath: A Curated Collection of Open Mathematics Problems, catalogue entry. Catalogue curation and metadata: CC BY 4.0. The source extract is retained; definitions, English summaries and qualifications are editorial additions.
\url{https://huggingface.co/datasets/ulamai/UnsolvedMath}.
\item[S2] Complexity Theory (2024), Oberwolfach Reports, DOI 10.4171/owr/2024/27, pp. 1517--1562
\url{https://doi.org/10.4171/owr/2024/27}.
\item[S3] Feng, Guo, Yin and Zhang, colorings
\url{https://arxiv.org/abs/2007.08091}.
\item Source-specific Oberwolfach Reports licences, including CC BY-SA 4.0 where applicable, are retained; see THIRD\_PARTY\_NOTICES.md.
\end{itemize}
\paragraph{Preparation and review.}
The original source and UnsolvedMath contributors retain attribution for the source extract. Codex (GPT-6; exact serving revision not exposed) assisted with definitions, English summaries, scope notes and formatting. No new human mathematical authorship or independently verified solution is asserted. Catalogue identity checked against the complete supplied registry and current repository. Distinct upstream records, including repeated questions and broad or damaged extracts, are retained for complete coverage of the source CS category. A repeated record is not a new independent conjecture. Complete mathematical-equivalence and current-literature review has not been performed. Preparation and final review follow contribution prompts C and D, with source gaps disclosed. Structural and rendering checks do not establish a new theorem or open status.
\label{end:statement}

\section{Research attempt}
\subsection{Literature, source context and scope}
The full original catalogue body and source qualifications are preserved above. The relevant primary workshop passages were read in the cached publisher report. Historical claims and subsequent attributed results are distinguished from the mathematical work below. A restricted example, standard lemma or known theorem is not counted as new progress on the entire original target. Missing assumptions and repeated or rhetorical extracts are not silently turned into different solved problems.
Codex (GPT-6; exact serving revision not exposed) supplied the auxiliary proof route used here. OpenAI \texttt{gpt-6-astra}, reasoning \texttt{ultra}, independently worked through its mathematical steps, wrote this exposition, and performed the finite checks. This is model collaboration, not independent human verification.

\subsection{A high-color coupling and the exact gap to the conjecture}
The actual primary conjecture is $O(1)$ spectral independence for all proper $q$-colorings when $q\ge\Delta+2$. We prove only a familiar rapid-mixing bound in the more restrictive regime $q>2\Delta$. A mixing estimate alone is not asserted to prove the source spectral-independence statement. Let $G$ be a finite simple graph with $n\ge1$ vertices, maximum degree $\Delta$, and integer $q\ge2$ with $q>2\Delta$.
At each step choose a vertex uniformly and recolor it uniformly from the colors absent at its neighbors, allowing its old color. Let $\pi$ be uniform on all proper colorings; this set is nonempty by sequential greedy coloring, since $q\ge\Delta+1$.
\begin{proposition}
Put $\kappa=(q-2\Delta)/(q-\Delta)>0$. From every proper starting coloring $\sigma$,
\[
\|P^t(\sigma,\cdot)-\pi\|_{\rm TV}\le n e^{-\kappa t/n}.
\]
\end{proposition}
\begin{proof}
Two colorings differing only at vertex $v$ have the same neighbor colors there, hence the same available set and identical forward/reverse transition probabilities. Other distinct states have zero transition in either direction. This proves detailed balance with the uniform law and therefore stationarity.
Couple any two chains by choosing the same vertex $v$. Let $A,B$ be their available color sets. Couple their uniform choices maximally: match each common color with probability $\min\{1/|A|,1/|B|\}$, then couple the remaining disjoint residual masses. Their disagreement probability is
\[
1-\frac{|A\cap B|}{\max\{|A|,|B|\}}.
\]
Let $d_v$ be the number of neighbors on which the colorings disagree. Each color in either directional difference of available sets must be supplied by a disagreeing neighbor in the other coloring. Thus both directional differences have size at most $d_v$, and $|A|,|B|\ge q-\Delta$. The disagreement probability is at most $d_v/(q-\Delta)$.
If $D$ denotes their current Hamming distance, updating $v$ removes its old disagreement and adds one with at most this probability. Averaging over $v$ gives
\[
\mathbb E[D_{t+1}\mid\sigma_t,\tau_t]
\le D_t-D_t/n+\frac{\sum_vd_v}{n(q-\Delta)}
\le(1-\kappa/n)D_t,
\]
since each disagreeing vertex is counted in at most $\Delta$ neighbor lists. Iteration gives $\mathbb ED_t\le n(1-\kappa/n)^t\le ne^{-\kappa t/n}$. Start the second chain at stationarity. For any event, the difference of the two probabilities is at most the probability of unequal coupled states, which is at most $\mathbb ED_t$. Taking the supremum over events proves the result.
\end{proof}
This direct all-pairs coupling needs no unproved path-coupling lemma. At the target boundary $q=\Delta+2$, its coefficient $\kappa$ is nonpositive once $\Delta\ge2$ (zero at two), so this estimate supplies no contraction there. Better influence control is precisely the missing step. Frozen colorings of $K_q$ at $q=\Delta+1$ do not refute the source conjecture, because they lie outside its range.
\subsection{Verification and remaining gaps}
Available-set total variations and the full Hamming contraction were exhaustively checked for small graphs and every pair of proper colorings. Transition-matrix stationarity and finite-time TV bounds were checked separately. The result is a known high-color regime, explicitly credited as an auxiliary route. Neither spectral independence at $q\ge\Delta+2$ nor arbitrary pinned-list variants are proved.

\section{References and source attribution}
Primary workshop report: \url{https://doi.org/10.4171/owr/2024/27}. Relevant passages read: 2024/27 p.1533, exact q>=Delta+2 spectral-independence conjecture and stated special regimes.

\url{https://arxiv.org/abs/2007.08058}. Read primary abstract on spectral independence for triangle-free colorings; its restricted parameter regime is not identified with the full source conjecture.

All retrievals were checked on 11 October 2026. Original question attribution remains in the preserved Sources section. External theorems are attributed and not claimed as new discoveries or fully verified merely by citation.


\paragraph{Integration review.} Codex (GPT-6; exact serving revision not exposed) reviewed the written arguments and expanded layout and source-reference wrappers. Repository metadata and the production web renderer were checked using Python 3.13 and JavaScript with MathJax 3.2.2. This model-assisted review is not independent human proof verification; no PDF compilation or formal proof checking is asserted.

\section{Completion Estimate}
\noindent\textbf{COMPLETION: 0\%}
No substantiated new progress toward the entire original target is claimed. The worked results and failed route are delimited above; a known result, illustrative model, or auxiliary obstruction is not counted as original-target completion.

\section{Final status and attribution}
\textbf{Attempt status: unresolved.} This describes this research attempt, not a new ruling on the source catalogue's status or on external literature claims.
The mathematics and this self-audit were generated by OpenAI \texttt{gpt-6-astra}, reasoning \texttt{ultra}, on 11 October 2026. Codex (GPT-6; exact serving revision not exposed) orchestrated the batch, read and reviewed its written mathematical arguments, and contributed formatting/repository review. No human mathematical authorship, independent human proof checking, or formal verification is claimed. The historical preparation paragraphs in the reproduced catalogue body concern that body only.

Contribution prompt A was applied to the complete definition and supplied source record, followed by the adversarial checks of prompt B. Local reproduction files are \texttt{batch025/checks.py}, \texttt{computations.json}, and the per-ID audit JSON. They were not public links at preparation time. Numerical checks are finite and do not prove unrestricted results. No LaTeX engine was found locally; compilation was not performed by this research agent.

\end{document}
